% Library unit: Bähr's rule. Printed once, before the first position that uses it
% (\xref{Section}{bahr-rule} or \uses{bahr-rule}); left out of a book that never uses it.

\paragraph{Bähr's rule}\label{bahr-rule}\index{King and pawn endings!Bähr's rule}
Two rook's pawns often block each other while one side has a passed pawn on the other wing. The
stronger side's king goes after the enemy rook's pawn; the defender takes the passed pawn and runs
back to the corner in front of his own rook's pawn. Who comes first? Bähr gave a rule for this in
1936; Dvoretsky offers a handier version~\cite{Dvoretsky:2008-endgame-manual}. Here it is with
White as the stronger side, the rook's pawns on the a-file and the passed pawn on the king side.
\begin{itemize}
\item If White's rook's pawn has crossed the middle of the board, he always wins.
\item The \emph{normal} position is a draw. In it, the rook's pawns stand on either side of the
middle (a4 against a5), and the passed pawn has reached or crossed the key diagonal h3--c8, so the
black king reaches c8 in time.
\item Every square the passed pawn stands behind this diagonal is an extra tempo for White, and so
is his king standing in front of the passed pawn instead of beside it. Every square the rook's
pawns stand further back than in the normal position is an extra tempo for Black.
\item White wins only if he has more tempi than Black.
\end{itemize}

\keysquares{8/8/5k2/8/5K2/p4P2/P7/8 w - - 0 1}{c8,d7,e6,f5,g4,h3}{White to move wins. The key diagonal h3--c8 is marked.}
Here is Dvoretsky's example~\cite{Dvoretsky:2008-endgame-manual}. The queenside pawns stand two
squares behind the normal position (a2 and a3 instead of a4 and a5): two tempi for Black. The
f3-pawn is two squares behind the key diagonal, which crosses the f-file on f5: two tempi for
White. His king stands in front of the pawn, not beside it: a third tempo. White wins 3:2:
1.\K e4! \K e6 2.\K d4 \K f5 3.\K c3 \K f4 4.\K b3 \K xf3 5.\K xa3 \K e4 6.\K b4 \K d5 7.\K b5
\K d6 8.\K b6 \K d7 9.\K b7 \K d6 10.a4, and the pawn queens.

1.\K e3? would be a terrible mistake. The white king is then beside the pawn, no longer in front
of it, so the count is only 2:2, and that is a draw: 1...\K e5! 2.\K d3 \K f4 3.\K c3 \K xf3
4.\K b3 \K e4 5.\K xa3 \K d5 6.\K b4 \K c6, and the black king is in time.
